% ===========================================================================
%  Homework 03 — Homework 03
%
%  DIVISION OF LABOUR (see AI_INSTRUCTIONS.md §7):
%    The assistant transcribes each assigned problem statement faithfully from
%    the assignment sheet in this folder's assignment/ directory, and emits an
%    empty, marked solution region beneath it.
%    The user types the mathematics inside those regions. No assistant writes
%    inside a SOLUTION region in normal mode.
% ===========================================================================
\documentclass[11pt]{article}
\usepackage{coursemacros}

\title{Homework 03}
\author{Mikey Ferguson}
\date{\today}

\begin{document}
\maketitle

% ---------------------------------------------------------------------------
% Assistant: replace this block with one problem/solution pair per assigned
% problem, in the order assigned, following the pattern below exactly.
% ---------------------------------------------------------------------------

\begin{problem}{1}
  (Edition 11, Chapter 3.) If $X$ and $Y$ are both discrete, show that
  \[
    \sum_{x} p_{X \mid Y}(x \mid y) = 1
  \]
  for all $y$ such that $p_Y(y) > 0$.
\end{problem}
% ===== SOLUTION 1 =====
Fix a value $y$ with $p_Y(y) > 0$. By the chain rule, $p(x,y) = p_X(x)\,p_{Y \mid X}(y \mid x)$,
which holds for every joint pmf and assumes nothing about independence. Hence
\[
  \sum_{x} p_{X \mid Y}(x \mid y)
    = \sum_{x} \frac{p_X(x)\,p_{Y \mid X}(y \mid x)}{p_Y(y)}
    = \frac{1}{p_Y(y)} \sum_{x} p_X(x)\,p_{Y \mid X}(y \mid x),
\]
where $1/p_Y(y)$ comes out of the sum because $y$ is held fixed and the sum runs over $x$ alone.
The remaining sum is the joint pmf marginalised over $x$, which is $p_Y(y)$:
\[
  \sum_{x} p_X(x)\,p_{Y \mid X}(y \mid x) = \sum_{x} p(x,y) = p_Y(y).
\]
Therefore
\[
  \sum_{x} p_{X \mid Y}(x \mid y) = \frac{p_Y(y)}{p_Y(y)} = 1. \qedsymbol
\]
% ===== END SOLUTION 1 =====

\begin{problem}{7}
  (Edition 11, Chapter 3.) Suppose $p(x,y,z)$, the joint probability mass
  function of the random variables $X$, $Y$ and $Z$, is given by
  \[
    \begin{array}{ll}
      p(1,1,1) = \tfrac{1}{8},  & p(2,1,1) = \tfrac{1}{4},\\[2pt]
      p(1,1,2) = \tfrac{1}{8},  & p(2,1,2) = \tfrac{3}{16},\\[2pt]
      p(1,2,1) = \tfrac{1}{16}, & p(2,2,1) = 0,\\[2pt]
      p(1,2,2) = 0,             & p(2,2,2) = \tfrac{1}{4}.
    \end{array}
  \]
  What is $\EEc{X}{Y = 2}$? What is $\EEc{X}{Y = 2, Z = 1}$?
\end{problem}
% ===== SOLUTION 7 =====
% TODO(mferguson): your work goes here.
% ===== END SOLUTION 7 =====

\begin{problem}{9 or 10}
  (Edition 11, Chapter 3. The sheet assigns either one of these.)

  \textbf{9.} Show in the discrete case that if $X$ and $Y$ are independent, then
  \[
    \EEc{X}{Y = y} = \EE{X} \qquad \text{for all } y.
  \]

  \textbf{10.} Suppose $X$ and $Y$ are independent continuous random variables.
  Show that
  \[
    \EEc{X}{Y = y} = \EE{X} \qquad \text{for all } y.
  \]
\end{problem}
% ===== SOLUTION 9 or 10 =====
% TODO(mferguson): your work goes here.
% ===== END SOLUTION 9 or 10 =====

\begin{problem}{12}
  (Edition 11, Chapter 3.) The joint density of $X$ and $Y$ is given by
  \[
    f(x,y) = \frac{e^{-x/y} e^{-y}}{y}, \qquad 0 < x < \infty,\ \ 0 < y < \infty.
  \]
  Show $\EEc{X}{Y = y} = y$.
\end{problem}
% ===== SOLUTION 12 =====
% TODO(mferguson): your work goes here.
% ===== END SOLUTION 12 =====

\begin{problem}{22}
  (Edition 11, Chapter 3.) Suppose that independent trials, each of which is
  equally likely to have any of $m$ possible outcomes, are performed until the
  same outcome occurs $k$ consecutive times. If $N$ denotes the number of
  trials, show that
  \[
    \EE{N} = \frac{m^{k} - 1}{m - 1}.
  \]
\end{problem}
% ===== SOLUTION 22 =====
% TODO(mferguson): your work goes here.
% ===== END SOLUTION 22 =====

\begin{problem}{31}
  (Edition 11, Chapter 3.) Each element in a sequence of binary data is either
  $1$ with probability $p$ or $0$ with probability $1 - p$. A maximal
  subsequence of consecutive values having identical outcomes is called a
  \emph{run}. For instance, if the outcome sequence is $1,1,0,1,1,1,0$, the
  first run is of length $2$, the second is of length $1$, and the third is of
  length $3$.
  \begin{enumerate}
    \item[(a)] Find the expected length of the first run.
    \item[(b)] Find the expected length of the second run.
  \end{enumerate}
\end{problem}
% ===== SOLUTION 31 =====
% TODO(mferguson): your work goes here.
% ===== END SOLUTION 31 =====

\begin{problem}{37}
  (Edition 11, Chapter 3.) A manuscript is sent to a typing firm consisting of
  typists $A$, $B$ and $C$. If it is typed by $A$, then the number of errors
  made is a Poisson random variable with mean $2.6$; if typed by $B$, then the
  number of errors is a Poisson random variable with mean $3$; and if typed by
  $C$, then it is a Poisson random variable with mean $3.4$. Let $X$ denote the
  number of errors in the typed manuscript. Assume that each typist is equally
  likely to do the work.
  \begin{enumerate}
    \item[(a)] Find $\EE{X}$.
    \item[(b)] Find $\Var(X)$.
  \end{enumerate}
\end{problem}
% ===== SOLUTION 37 =====
% TODO(mferguson): your work goes here.
% ===== END SOLUTION 37 =====

\begin{problem}{40}
  (Edition 11, Chapter 3.) A prisoner is trapped in a cell containing three
  doors. The first door leads to a tunnel that returns him to his cell after two
  days of travel. The second leads to a tunnel that returns him to his cell
  after three days of travel. The third door leads immediately to freedom.
  \begin{enumerate}
    \item[(a)] Assuming that the prisoner will always select doors $1$, $2$ and
      $3$ with probabilities $0.5$, $0.3$, $0.2$, what is the expected number of
      days until he reaches freedom?
    \item[(b)] Assuming that the prisoner is always equally likely to choose
      among those doors that he has not used, what is the expected number of days
      until he reaches freedom? (In this version, for instance, if the prisoner
      initially tries door $1$, then when he returns to the cell, he will now
      select only from doors $2$ and $3$.)
    \item[(c)] For parts (a) and (b) find the variance of the number of days
      until the prisoner reaches freedom.
  \end{enumerate}
\end{problem}
% ===== SOLUTION 40 =====
% TODO(mferguson): your work goes here.
% ===== END SOLUTION 40 =====

\end{document}
